\mxnewpage\section*{EV2-T1 -- Blanc 1}\label{EV2-T1}
\addcontentsline{toc}{section}{EV2-T1 -- Blanc 1}
\textbf{Sujet d’entraînement}. Chapitres 1 à 7.\par Durée : 2 h. Sans calculatrice. Note sur 20. 14–16 janvier 2027.
\subsection*{Partie 1 -- Automatismes (6 points)}
Douze questions, 0{,}5 point chacune. Une seule réponse correcte ; aucune justification demandée. Une réponse incorrecte ou absente rapporte 0, sans retrait de points.
\begin{enumerate}
\item Calculer $2+2\times3$.\par A. $8$\quad B. $12$\quad C. $11$\quad D. $7$
\item Calculer $\frac12+\frac13$.\par A. $\frac{1}{6}$\quad B. $\frac{5}{3}$\quad C. $\frac{2}{5}$\quad D. $\frac{5}{6}$
\item Calculer $\dfrac{2^2}{2^{0}}$.\par A. $2$\quad B. $7$\quad C. $1$\quad D. $4$
\item Une valeur de 50 augmente de 20 \%. Valeur finale ?\par A. $40$\quad B. $60$\quad C. $70$\quad D. $10$
\item Une valeur de 80 est obtenue après une baisse de 20 \%. Valeur initiale ?\par A. $96$\quad B. $100$\quad C. $64$\quad D. $60$
\item Hausse de 25 \% puis baisse de 20 \% : coefficient global ?\par A. $\frac{4}{5}$\quad B. $\frac{9}{10}$\quad C. $1$\quad D. $\frac{21}{20}$
\item La solution de $2x+2=0$ vaut\par A. $-4$\quad B. $-1$\quad C. $1$\quad D. $4$
\item La pente entre $(0;2)$ et $(2;8)$ vaut\par A. $6$\quad B. $2$\quad C. $3$\quad D. $-3$
\item La moyenne de $2, 4, 6$ vaut\par A. $4$\quad B. $2$\quad C. $12$\quad D. $5$
\item Un secteur de $90^\circ$ représente quelle fraction du total ?\par A. $\frac{1}{3}$\quad B. $\frac{3}{4}$\quad C. $\frac{1}{4}$\quad D. $\frac{1}{2}$
\item Une probabilité vaut $\dfrac{2}{5}$. Son complément vaut\par A. $\frac{2}{5}$\quad B. $\frac{5}{2}$\quad C. $\frac{1}{5}$\quad D. $\frac{3}{5}$
\item Pour $f(x)=x^2+2$, $f(-2)$ vaut\par A. $4$\quad B. $8$\quad C. $6$\quad D. $-2$
\end{enumerate}
\mxnewpage\subsection*{Partie 2 -- Exercices indépendants (14 points)}
Toute démarche doit être justifiée. Les résultats exacts sont attendus sauf indication contraire.
\mxneed{12}
\subsection*{Exercice 1 -- Fabrication et contrôle (5 points)}
Deux ateliers A et B produisent respectivement 30 pour cent et 70 pour cent des objets d’une entreprise. Parmi les objets de A, 10 pour cent sont défectueux ; cette proportion est de 20 pour cent chez B. On choisit un objet au hasard dans la production totale. Les événements $A$ et $D$ signifient respectivement « provenir de l’atelier A » et « être défectueux ».
\begin{enumerate}
\item Construire un arbre pondéré dont le premier niveau distingue les deux ateliers. \hfill(1 pt)
\item Calculer $P(D)$, puis la probabilité qu’un objet défectueux provienne de A. \hfill(2 pts)
\item Les événements $A$ et $D$ sont-ils indépendants ? Justifier numériquement. \hfill(2 pts)
\end{enumerate}
\mxneed{12}
\subsection*{Exercice 2 -- Minimum d’un quotient (5 points)}
La fonction $f$ est définie sur l’intervalle $]0;+\infty[$ par $f(x)=x+\dfrac9x$. On étudie ses variations et certaines tangentes à sa courbe. Les coordonnées demandées sont à donner sous forme exacte.
\begin{enumerate}
\item Calculer et factoriser la dérivée. \hfill(1 pt)
\item Étudier les variations et donner le minimum. \hfill(2 pts)
\item Déterminer les points où la tangente a pour pente $\dfrac{3}{4}$. \hfill(2 pts)
\end{enumerate}
\mxneed{12}
\subsection*{Exercice 3 -- Projeté et cercle (4 points)}
Dans un repère orthonormé, on considère le point $P(5;2)$ et la droite $d$ d’équation $x+y-3=0$. On note $H$ le projeté orthogonal de $P$ sur $d$.
\begin{enumerate}
\item Donner un vecteur normal non nul à $d$. \hfill(1 pt)
\item Déterminer les coordonnées de $H$ et vérifier les conditions d’appartenance et d’orthogonalité. \hfill(1 pt)
\item Donner une équation du cercle de centre $P$ passant par $H$. \hfill(2 pts)
\end{enumerate}
